% =========================================================================
% 九、模型评价
% =========================================================================

\section{模型评价与改进（问题一至问题四）}

本节只陈述有证据支持的具体性质。评价分两层：先说清``好在哪、由什么证据支持、在什么范围内成立''，再说``局限对应哪一项实质性改进''。

\subsection{模型的具体优势及其证据}

\begin{enumerate}
  \item \textbf{子问题可求全局最优且计算量小。}问题一至问题四的计划层与调整层都是线性规划，规模为百至千量级变量，用 HiGHS 求解在秒级完成；问题一的充放电互斥约束经命题保证线性松弛在最优值意义下精确（\S\ref{sec:q1-relax}），问题三的含分段结算规则也被化归为\emph{单个}线性规划（式~\eqref{eq:q3-adjust}），无需二元变量。相应地，{\bfseries 主方案给出的最优性是针对给定输入、给定约束的确定性子问题的最优性}，不是对整条滚动策略的整体最优性。
  \item \textbf{执行规则完全因果、可行域明确。}式~\eqref{eq:q2-exec} 的尽限充放电规则只依赖当前区间的实测净负载与当前实际库存，充放互斥由净缺口符号自动决定，实际库存始终落在 $[1200,10800]$ kWh 内；把实测净负载设为计划所用序列时，执行器逐位复现线性规划的解（\S\ref{sec:q2-exec}）。
  \item \textbf{实际库存跨日连续。}$E_H^{(d)}$ 直接结转成为 $E_0^{(d+1)}$，评价期 334 天内跨日衔接误差为 0，不存在人为重置库存带来的虚假收益。
  \item \textbf{账单可独立复算。}Q1 的购电费、Q2 的计划费与紧急费、Q4 两分支的合同费与附加费，均由不调用调度代码的独立路径重算并逐位吻合（\S\ref{sec:q1-constraint}、\S\ref{sec:q2-exec}、\S\ref{sec:q4-check}）。
  \item \textbf{$2\times2$ 对照能区分``信息''与``调整机会''。}问题三用``是否用 0:00 预报通道 $\times$ 是否允许日内调整''的四格设计，把两类改进的贡献与二者 9.62 万元的重叠分开计量（\S\ref{sec:q3-matrix}），避免了单项消融无法处理的相互替代问题。
  \item \textbf{预测器形式与参数取值受因果协议约束。}问题二的 $(\beta,W,\gamma,N_{\mathrm{DOW}})$ 按逐月 walk-forward 选取，问题三的 $\kappa_g$ 只用 $d$ 之前 30 天逐日重估，问题四的日内价格通道经``扰动未来电价、预测不变''的检验。问题三、问题四的固定配置为同年探索性选定，本文如实标注而不宣称独立测试（\S\ref{sec:q4-limit}）。
  \item \textbf{结论的边界被显式写出。}问题三``日内新增预报未进一步降低费用''限于所比较的全融合与门控两种策略、以及本文的结算口径（\S\ref{sec:q3-other}）；问题四的价格侧差额限于 24 小时固定策略（\S\ref{sec:q4-ablation}）；容量边际价值限于``固定下限、抬升上限''的扰动路径（\S\ref{sec:q1-sens}）。
\end{enumerate}

\subsection{局限与对应的实质性改进}

\begin{table}[htbp]
  \centering
  \caption{主要局限与对应的实质性改进方向}
  \label{tab:limits}
  \small
  \begin{tabularx}{\textwidth}{XX}
    \toprule
    当前的局限 & 对应的实质改进 \\
    \midrule
    预测精度指标（MAE）与实际结算费用方向不一致：问题三、问题四都出现``预测更准但费用未降'' & 在历史滚动验证中\textbf{直接用真实结算费用}为目标，去校准风险裕量水平与预报融合权重，而不是先优化 MAE 再检验费用 \\
    日前规划没有把未来的调整机会内生化：计划层按点预测求解，未考虑 6/12/18 时仍可重定合同 & 研究受非预见性约束的多阶段决策，或保持同一因果执行器、以回放费用重评各候选计划 \\
    48 小时时域的次日输入为``当日预测平铺''的持久性预测 & 加入只依赖当时历史的星期与季节条件预测，并在同一窗口下作对照 \\
    固定策略按同年结果探索性选择（问题三的 $\kappa_g$ 网格、问题四的 $(\lambda_{\mathrm{p}},\delta)$ 与时域） & 在后续独立年份上验证，或按本文方式如实保留``探索性、未做外部测试''的定位 \\
    未计电池退化、未建模售电与逆变器限制 & 在推广模型中加入循环成本、售电价与逆变器约束，并重新检验互斥松弛的精确性条件 \\
    风险裕量按各时段独立取分位数，未建模时段相关性 & 引入时段间相关结构（如情景集合或分层分位），并检验其对紧急费用的影响 \\
    \bottomrule
  \end{tabularx}
\end{table}

\subsection{关于创新点的表述}

本文的创新应表述为\textbf{针对本题所作的结构设计及其验证}，而不宣称算法本身原创：其一是把含双向违约价的调整决策化归为单个线性规划（式~\eqref{eq:q3-adjust}），使合同重定仍可求子问题最优；其二是用 $2\times2$ 对照分离``新增信息''与``调整机会''的贡献并量化二者重叠；其三是在波动电价下先厘清价格缩放不变性的适用边界，再以加权裕量与 48 小时时域分别处理价格结构与库存配置，并用同政策对照与放松下界界定各结论的范围。本文的主方案是线性规划与尽限反馈的组合，故不以动态规划的状态价值函数或深度学习方法作为全篇核心贡献；建模过程中测试过但未取得优势的做法（如逐情景回放择优的在线选择器）只作为负结果如实报告，不与主模型并列。
